% concatenated from SimonConjecture.tex
\documentclass[11pt,reqno]{amsart}
\usepackage{amsthm,amsfonts,amsmath,amssymb}
\usepackage[pagebackref=true,backref=true,bookmarksopen=true]{hyperref}
\usepackage[letterpaper,margin=1in]{geometry}

\newtheorem{theorem}{Theorem}
\newtheorem{cnj}{Conjecture}
\newtheorem{lemma}{Lemma}

\newcommand{\R}{{\mathbb R}}

\begin{document}

\title{Separation for Schr\"odinger operators. A counterexample to Simon's conjecture}
\author{M.A. Perelmuter}
\address{SCAD Soft Ltd., 3a Osvity Street, Kyiv, 03037, Ukraine}
\email{mikeperelmuter@gmail.com}
\keywords{Separation property, Schr\"odinger operator, counterexample}
\subjclass[2020]{35J10 35P05}

\begin{abstract}
We construct, for every $d\ge 5$, a nonnegative radial potential
$V\in L^\infty_{\mathrm{loc}}(\mathbb{R}^d\setminus\{0\})\cap L^1(\mathbb{R}^d)$ for which
$V(-\Delta\dotplus V)^{-1}$ is not bounded on $L^2(\R^d)$.
\end{abstract}

\maketitle
\section{Introduction}
Let $H_0=-\Delta$ on $L^2(\R^d)$. $H=H_0\dotplus V$ is defined as a self-adjoint operator via the form sum.
Clearly, $\mathcal{D}(H_0)\cap \mathcal{D}(V) \subset \mathcal{D}(H)$ and any condition implying
\begin{equation}\label{Separate}
\mathcal{D}(H) \subset \mathcal{D}(V)
\end{equation}
implies $H=H_0+V$. W.~N.~Everitt and M.~Giertz, in \cite{Ever}, coined the term {\it separation property} for relation \eqref{Separate}.
There are many publications devoted to sufficient conditions for the separation property of various differential operators. We mention only a few of them here: \cite{Ever}, \cite{Ever73}, \cite{Ever77}, \cite{Davies}, \cite{Davies2}, \cite{Otelbaev}, \cite{Kalf}, \cite{Okazawa},  \cite{GS}, \cite{MS}.

In 1983 B.~Simon \cite{Simon2} formulated the following conjecture
\begin{cnj}\label{SimonCnj}
Let $V(x)=V(|x|)$ on $\R^d; d\geq 5$. Suppose
\begin{enumerate}
  \item $V\geq 0$
  \item $V$ is bounded on any closed set avoiding $0$.
\end{enumerate}
Then $H_0+V$ is closed on $\mathcal{D}(H_0)\cap \mathcal{D}(V)$.
\end{cnj}
The fact that the conjecture is false for $d=4$ was proved by B.~Simon \cite{Simon2} himself. To the best of our knowledge, this conjecture has remained open for $d\geq 5$, and the purpose of this note is to construct a counterexample.
\par
We construct a radial potential $V\ge0$, singular only at the origin, for which
$V(-\Delta+V)^{-1}$ fails to be bounded on $L^2(\R^d)$. The potential
is supported on a countable family of thin, pairwise disjoint spherical shells
$A_n=[r_n,\,r_n+w_n],\ r_n \to 0$, accumulating at the origin, on which $V$ takes the value
$h_n\sim r_n^{\,1-d}$. The widths $w_n$ are chosen so that the quantity $h_n w_n$, which may be viewed as the mass of each shell, is independent of $n$: the shells become thinner and taller as $n\to\infty$.
\par
For this potential we construct $f\in L^2(\R^d)$ and $u\in \mathcal{D}(H)$ such that $Hu=f$ and $Vu\notin L^2(\R^d)$. By spherical symmetry the construction reduces to a one-dimensional Bessel-type equation, which is then converted into a Volterra integral equation with a special integral kernel.

\section{A counterexample}
\begin{theorem}
Let $d\ge 5$. There exists a radial potential
$$0\le V\in L^\infty_{\mathrm{loc}}(\R^d\setminus\{0\})\cap L^1({\R^d})$$
such that the operator $V(-\Delta\dotplus V)^{-1}$ is not bounded on $L^2(\R^d)$.
\end{theorem}
\begin{proof}
Let $r_n=2^{-n},\ w_n=r_n^{d-1},\ n\ge 2,$ and $A_n=[r_n,r_n+w_n]$.
Since $\frac{w_n}{r_n}=r_n^{d-2}<1,$ we have $r_n+w_n<2r_n=r_{n-1},$ so the intervals $A_n$ are pairwise disjoint.
\par
Fix $\lambda>0$ and define $h_n=\frac{\lambda}{w_n}=\lambda r_n^{1-d}$,
$$V(r)=
\begin{cases}
h_n,&r\in A_n,\\
0,&r\notin\bigcup_{n\ge 2}A_n.
\end{cases}$$
Thus $V$ is radial and nonnegative. Since every compact subset of $\R^d\setminus\{0\}$ intersects only finitely many layers, we have $V\in L^\infty_{\mathrm{loc}}(\R^d\setminus\{0\}).$ The inclusion $V\in L^1(\R^d)$ is a consequence of the equality $V(r)=0$ for $r\notin\bigcup_{n\ge 2}A_n$ and the estimate
\begin{equation*}
\|V\|_{L^1(\R^d)}=\omega_{d-1}\int\limits_0^1 r^{d-1}V(r)\,dr\le \omega_{d-1} 2^{d-1}\lambda\sum_{n\ge 2}r_n^{d-1}<\infty,
\end{equation*}
where $\omega_{d-1}$ is the surface area of the unit $(d-1)$-dimensional sphere.
\par
Define the quadratic form
$$q[u]=\int\limits_{\R^d}|\nabla u|^2\,dx+\int\limits_{\R^d}V|u|^2\,dx$$
 with domain $\mathcal{D}(q)=H^1(\R^d)\cap L^2(\R^d,V\,dx)$. Since $0\le V\in L^1(\R^d)$, this is a densely defined closed nonnegative quadratic form. Let $H$ be the associated self-adjoint operator (the form realization of $-\Delta+V$).
 Let $f(x)=\chi_{\{1\le |x|\le 2\}}\in L^2(\R^d)$ (the indicator function of the set $\{1\le |x|\le 2\}$). We shall construct a radial solution $u\in \mathcal{D}(H)$ such that $Hu=f$. It will then follow that $u=H^{-1}f=(-\Delta\dotplus V)^{-1}f.$
\par
To construct the function $u$, we will consider three intervals: $(0,1), [1,2], (2,\infty)$.
\par
1. As a first step we consider a solution near the origin ($0 < r < 1$).
By spherical symmetry, the construction of the function $u$ can be reduced to solving a one-dimensional Bessel-type equation.
Consider the radial homogeneous equation
$$-U''(r)-\frac{d-1}{r}U'(r)+V(r)U(r)=0.$$
Define $U$ as the solution of the Volterra equation
\begin{equation}\label{UVoltera}
U(r)=1+\frac{1}{d-2}\int\limits_0^r\left(s-r^{2-d}s^{d-1}\right)V(s)U(s)\,ds.
\end{equation}
Since $V\ge 0$, the kernel is nonnegative, and $0\le s-r^{2-d}s^{d-1}\le s,\  0<s<r$. Moreover,
\begin{equation}\label{rV}
\int\limits_0^1 rV(r)\,dr=\sum_{n\ge 2}h_n\int\limits_{r_n}^{r_n+w_n}r\,dr\le\sum_{n\ge 2}h_n(r_n+w_n)w_n\le
2\lambda\sum_{n\ge 2}r_n<\infty.
\end{equation}
By the standard existence and uniqueness theory for linear Volterra integral equations of the second kind with integrable kernels, equation \eqref{UVoltera} has a unique continuous solution on $[0,1]$.
\par
Differentiating \eqref{UVoltera}, we obtain
$$U'(r)=r^{1-d}\int\limits_0^r s^{d-1}V(s)U(s)\,ds\ge 0,$$
and therefore
\begin{equation}\label{Ule1}
U(r)\ge U(0)=1.
\end{equation}
Moreover, \eqref{UVoltera} gives
$$U(r)\le 1+\frac{1}{d-2}\int\limits_0^r sV(s)U(s)\,ds.$$
Hence Gronwall's inequality and \eqref{rV} imply
\begin{equation}\label{Uge1}
1\le U(r)\le\exp\left(\frac{1}{d-2}\int\limits_0^1sV(s)\,ds\right)=:C_0<\infty.
\end{equation}
We next record an estimate for $U'$ valid on all of $(0,1)$, not merely on the layers $A_n$ themselves. For $r\in(0,1)$ let $n=n(r)\ge 2$ be the unique integer with $r\in[r_n,r_{n-1})$, where $r_1:=1$.
Since $A_k\subset[r_k,r_{k-1})$ for every $k\ge 2$, the set $\{s\le r:V(s)\ne 0\}$ is contained in $\bigcup_{k\ge n}A_k$, and $r^{1-d}\le r_n^{1-d}$ because $r\ge r_n$. Hence
\begin{equation}\label{UPrime}
\begin{aligned}
U'(r)
&\le
C_0r^{1-d}\int\limits_0^r s^{d-1}V(s)\,ds\\
&\le
C_0 r_n^{1-d}
\sum_{k\ge n}h_kw_k (r_k+w_k)^{d-1}\\
&\le
C\lambda r_n^{1-d}\sum_{k\ge n}r_k^{d-1}
\le C.
\end{aligned}
\end{equation}
Here and below $C$ denotes a finite constant independent of $n$ (hence of $r$). Since this bound holds for every $r\in(0,1)$, we conclude
$$\int\limits_0^1|U'(r)|^2r^{d-1}\,dr\le C^2\int\limits_0^1r^{d-1}\,dr<\infty.$$
\par\medskip
On $(0,1)$ put
\begin{equation}\label{ucU}
u(r)=cU(r),
\end{equation}
where $c>0$ will be chosen below so that the solutions can be matched continuously with their first derivatives at the interval endpoints.
\par\medskip
2. On $1\le r \le 2$, $V=0$ and $f=1$. Using the standard representation of the Laplace operator acting on radial functions, we have
\begin{equation}\label{r12}
-\left(r^{d-1}u'(r)\right)'=r^{d-1}.
\end{equation}
With the initial conditions $u(1)=cU(1)$ and $u'(1)=cU'(1)$, there exists a unique solution on $[1,2]$.
\par\medskip
3. For $r>2$, we have $V=f=0$.  We consider the homogeneous equation
$$-\bigl(r^{d-1}v'(r)\bigr)'=0$$
and choose the decaying solution
$$v(r)=b(c)r^{2-d},$$
where $b(c)$ is chosen so that the solution is continuous at $r=2$, namely, $b(c)=2^{d-2}u(2)$. For this choice, we have $v'(2)=-\frac{d-2}{2}u(2)$.
Consequently, in order for the solution defined on $[1,2]$ to be $C^1$-matched with the exterior solution at $r=2$, we need $u(2)+\frac{2}{d-2}u'(2)=0$.
\par
\par
By Lemma~\ref{lem:gluing} below, the unique solution of \eqref{r12} on $[1,2]$ satisfying $u(1)=cU(1)$, $u'(1)=cU'(1)$ satisfies
$$u(2)+\frac{2}{d-2}u'(2)=c\left(U(1)+\frac{U'(1)}{d-2}\right)-\frac{3}{2(d-2)}.$$
%The general solution of \eqref{r12} is $u(r)=A+Br^{2-d}-\dfrac{r^2}{2d}$, where $-\dfrac{r^2}{2d}$ is a particular solution of \eqref{r12} and $A+Br^{2-d}$ is the general solution of the homogeneous equation. The conditions $u(1)=cU(1)$ and $u'(1)=cU'(1)$ determine $A$ and $B$ in terms of $c,\ U(1),\ U'(1)$; substituting into $u(2)+\frac{2}{d-2}u'(2)$ and simplifying, we obtain
%$$u(2)+\frac{2}{d-2}u'(2)=c\left(U(1)+\frac{U'(1)}{d-2}\right)-\frac{3}{2(d-2)}.$$
Hence, there exists a unique value of $c>0$ for which $u(2)+\frac{2}{d-2}u'(2)=0$.
For this choice of $c$, defining
$$u(r)=b(c)r^{2-d},\qquad r\geq 2,$$
gives a $C^1$ solution across $r=2$.
\par\medskip
Since $d\ge 5$,
$$\int\limits_2^\infty |u(r)|^2r^{d-1}\,dr=b^2\int\limits_2^\infty r^{3-d}\,dr<\infty.$$
Together with \eqref{Uge1}, this gives $u\in L^2(\R^d)$.
\par\medskip
Now our function $u$ is defined for all $r\in (0,\infty)$ and we have to verify that $u$ belongs to $\mathcal{D}(H)$. We first prove that $u\in H^1(\R^d)$. Near the origin, \eqref{UPrime} implies
$$\int\limits_0^1|U'(r)|^2r^{d-1}\,dr<\infty.$$
Since $u=cU$ on $(0,1)$, the same integrability property holds for $u$.
On $(1,2)$ the solution is smooth, and for $r\ge 2$ the formula $u(r)=br^{2-d}$ gives $|u'(r)|\lesssim r^{1-d}$. Hence
$$\int\limits_2^\infty|u'(r)|^2r^{d-1}\,dr<\infty.$$
Thus $u\in H^1(\R^d)$.
\par
Next we estimate the potential part of the quadratic form. By \eqref{Uge1}, $u(r)\le cC_0,\  0<r\le 1$.
Therefore, using $V\in L^1(\R^d)$ and $V=0$ on $1\le r$, we have
$$\int\limits_{\R^d}V|u|^2\,dx =\int\limits_{|x|<1}V|u|^2\,dx\le c^2C_0^2\|V\|_{L^1(\R^d)}<\infty.$$
Thus $u\in \mathcal{D}(q)$, which yields a global radial solution of
$-\Delta u+Vu=f$. This differential equation holds classically away from the layer endpoints, and $u$ and $u'$ are continuous across all these endpoints. Hence it holds in the distributional sense on $\R^d$:
\begin{equation}\label{Distributional}
-\Delta u+Vu=f.
\end{equation}
Here $Vu\in L^1_{\mathrm{loc}}(\R^d)$ because $u$ is locally bounded and $V\in L^1(\R^d)$.
\par
Hence, for every $\varphi\in C_c^\infty(\R^d)$,
$$q[u,\varphi]=\int_{\R^d}f\varphi\,dx.$$
Since $C_c^\infty(\mathbb R^d)$ is a form core for $q$ (see, for example, \cite[Theorem 2.1]{Simon1}), this identity extends to every $\varphi\in\mathcal D(q)$. By the definition of the operator associated with the closed quadratic form $q$, it follows that
$$u\in\mathcal D(H),\quad Hu=f.$$
\par\medskip
To complete the proof of the theorem, it remains to show that $Vu\notin L^2(\R^d)$.
\par
By \eqref{Ule1} and \eqref{ucU}, $u(r)\ge c,\  0<r<1$. On $A_n$ we have $V=h_n$. Hence
$$\int\limits_{A_n}(Vu)^2\,dx = \omega_{d-1}h_n^2 \int\limits_{r_n}^{r_n+w_n}u(r)^2r^{d-1}\,dr\ge c^2\omega_{d-1}h_n^2r_n^{d-1}w_n.$$
Since $h_n=\frac{\lambda}{w_n},\  w_n=r_n^{d-1}$, we obtain $h_n^2r_n^{d-1}w_n=\lambda^2$.
Consequently
$$\int\limits_{A_n}(Vu)^2\,dx\ge c^2\omega_{d-1}\lambda^2=:\delta>0,$$
where $\delta$ is $n$-independent.
\par
The sets $A_n$ are pairwise disjoint, and therefore
$$\|Vu\|_{L^2(\R^d)}^2\ge\sum_{n\ge 2}\int\limits_{A_n}(Vu)^2\,dx\ge\sum_{n\ge 2}\delta=\infty.$$
Thus $Vu\notin L^2(\R^d)$. Hence $V(-\Delta\dotplus V)^{-1}$ is not a bounded operator on $L^2(\R^d)$.
\end{proof}

\section{Appendix: the gluing identity}
\begin{lemma}\label{lem:gluing}
Let $d\ge 3$, and let $u$ be the solution of
$$-\left(r^{d-1}u'(r)\right)'=r^{d-1},\qquad 1\le r\le 2,$$
subject to the initial conditions $u(1)=cU(1)$, $u'(1)=cU'(1)$, where $c,\,U(1),\,U'(1)$ are given real constants. Then
$$u(2)+\frac{2}{d-2}u'(2)=c\left(U(1)+\frac{U'(1)}{d-2}\right)-\frac{3}{2(d-2)}.$$
\end{lemma}
\begin{proof}
The general solution of the inhomogeneous equation is
$$u(r)=A+Br^{2-d}-\frac{r^2}{2d} \quad \text{and} \quad u'(r)=B(2-d)r^{1-d}-\frac{r}{d}.$$
Imposing the initial conditions at $r=1$:
$$u(1)=A+B-\frac{1}{2d}=cU(1),\qquad u'(1)=B(2-d)-\frac1d=cU'(1).$$
These equations give
$$B=-\frac{cU'(1)+\frac1d}{d-2}=-\frac{cU'(1)}{d-2}-\frac{1}{d(d-2)},\quad A=cU(1)-B+\frac{1}{2d}.$$
Now evaluate at $r=2$:
$$u(2)=A+B\,2^{2-d}-\frac2d,\qquad u'(2)=B(2-d)2^{1-d}-\frac2d.$$
The second equality gives
$$\frac{2}{d-2}u'(2)=\frac{2}{d-2}\Big[B(2-d)2^{1-d}-\frac2d\Big]=-B\,2^{2-d}-\frac{4}{d(d-2)}.$$
Adding this to $u(2)$:
$$u(2)+\frac{2}{d-2}u'(2)=A-\frac2d-\frac{4}{d(d-2)}.$$
Substituting $A=cU(1)-B+\dfrac{1}{2d}$ and then $B=-\dfrac{cU'(1)}{d-2}-\dfrac{1}{d(d-2)}$:
$$u(2)+\frac{2}{d-2}u'(2)=cU(1)+\frac{cU'(1)}{d-2}+\frac{1}{d(d-2)}-\frac{3}{2d}-\frac{4}{d(d-2)}=c\left(U(1)+\frac{U'(1)}{d-2}\right)-\frac{3}{d(d-2)}-\frac{3}{2d}.$$
Finally, the remaining constants combine as
$$\frac{3}{d(d-2)}+\frac{3}{2d}=\frac{3}{2(d-2)},$$
which yields the claimed identity.
\end{proof}


\smallskip
\noindent\textbf{Acknowledgements}\quad
The author acknowledges the use of AI tools. All mathematical arguments and proofs in the final manuscript were checked and written by the author.



\begin{thebibliography}{99}

\bibitem{Davies} Davies~E.~B.:
\newblock Some norm bounds and quadratic form inequalities for Schr\"odinger operators.
\newblock J. Oper. Theory, 1983, \textbf{9}, 147--162.

\bibitem{Davies2} Davies~E.~B.:
\newblock Some norm bounds and quadratic form inequalities for Schr\"odinger operators. II.
\newblock J. Oper. Theory, 1984, \textbf{12}, 177--196.

\bibitem{Ever} Everitt~W.~N., Giertz~M.:
\newblock Some properties of the domains of certain differential operators.
\newblock Proc. London Math. Soc. III, 1971, \textbf{23}, 301--324.

\bibitem{Ever73} Everitt~W.~N., Giertz~M.:
\newblock An Example concerning the Separation Property of Differential Operators.
\newblock Proc. R. Soc. Edinb., Sect. A, 1973, \textbf{71} (1971--1973), 159--165.

\bibitem{Ever77} Everitt~W.~N., Giertz~M.:
\newblock Inequalities and separation for Schr\"odinger type operators in $L^2(\R^n)$.
\newblock Proc. R. Soc. Edinb., Sect. A, 1977, \textbf{79}, No. 3-4, 257--265.

\bibitem{GS} Goldstein~J.~A., Svirsky~R.:
\newblock On a Domain Characterization of Schr\"odinger Operators with Gradient Magnetic Vector Potentials and Singular Potentials.
\newblock Proc. Amer. Math. Soc., 1989, \textbf{105}, No. 2, 317--323.

\bibitem{Kalf} Kalf~H.:
\newblock A note on the domain characterization of certain Schr\"odinger operators with strongly singular potentials.
\newblock Proc. R. Soc. Edinb., Sect. A, 1984, \textbf{97}, 125--130.

\bibitem{MS} Metafune~G., Schnaubelt~R.:
\newblock The domain of the Schr\"odinger operator $-\Delta+x^2y^2$.
\newblock Note Mat., 2006, \textbf{25}, No. 2, 97--103.

\bibitem{Okazawa} Okazawa~N.:
\newblock An $L^p$ theory for Schr\"odinger operators with nonnegative potentials.
\newblock J. Math. Soc. Japan, 1984, \textbf{36}, 675--688.

\bibitem{Otelbaev} Otelbaev~M.~O.:
\newblock Coercive estimates and separability theorems for elliptic equations in ${\mathbf R}^n$.
\newblock Proc. Steklov Inst. Math., 1983, \textbf{161}, 195--217.

\bibitem{Simon1} Simon~B.:
\newblock Maximal and minimal Schr\"odinger forms.
\newblock J. Oper. Theory, 1979, \textbf{1}, 37--47.

\bibitem{Simon2} Simon~B.:
\newblock Hardy and Rellich inequalities in non-integral dimension.
\newblock J. Oper. Theory, 1983, \textbf{9}, 143--146.

\end{thebibliography}
\end{document}
